\section{ТЕОРЕТИЧНІ ЗАСАДИ ФОРМУВАННЯ ІНФОРМАЦІЙНО-КОМУНІКАТИВНОЇ КОМПЕТЕНТНОСТІ МАЙБУТНІХ ДИЗАЙНЕРІВ}\label{sec:1}

У розділі висвітлено концептуальні основи формування інформаційно-комунікативної компетентності (ІКК) у системі професійної підготовки майбутніх дизайнерів. Розкрито сутність, структуру та специфіку ІКК як педагогічної категорії, що відповідає сучасним вимогам дизайнерської діяльності в умовах цифрової трансформації та стрімкого впровадження штучного інтелекту. Здійснено термінологічний аналіз суміжних понять (ІКТ-компетентність, цифрова компетентність, медіакомпетентність, інформаційна грамотність), що уможливив уточнення категоріального апарату дослідження. Обґрунтовано авторську 4-компонентну модель ІКК дизайнера, адаптовану до реалій професійної діяльності в еру генеративного штучного інтелекту. Проаналізовано трансформацію дизайн-освіти під впливом ШІ-технологій, стан розвитку компетентностей у контексті європейської інтеграції, а також синтезовано інтегративну теоретичну рамку дослідження.

\subsection{ІКК як педагогічна категорія: поняття, структура, специфіка дизайнера}\label{subsec:1.1}

Формування інформаційно-комунікативної компетентності (ІКК) майбутніх дизайнерів є актуальною педагогічною проблемою, зумовленою радикальною трансформацією дизайнерської професії під впливом цифрових технологій та штучного інтелекту. Сучасний дизайнер діє в інформаційно насиченому середовищі, де пошук, оброблення та візуалізація професійно значущої інформації, цифрова комунікація та критичне оцінювання технологічних інновацій набувають статусу ключових професійних дій. У цьому контексті теоретичне осмислення ІКК як педагогічної категорії потребує системного аналізу суміжних понять, визначення структурних компонентів та розкриття специфіки ІКК саме для дизайнерської професії.

Проблема формування компетентностей, пов'язаних з інформацією, комунікацією та технологіями, має тривалу історію в педагогічній науці. Починаючи з 1990-х років, у вітчизняному та зарубіжному дискурсі паралельно розвивалися кілька термінологічних ліній, що відображали різні аспекти взаємодії людини з інформаційним середовищем: інформаційна грамотність (information literacy), ІКТ-компетентність (ICT competence), цифрова компетентність (digital competence), медіакомпетентність (media competence) та інформаційно-комунікативна компетентність (ІКК). Неоднозначність трактувань ускладнює наукову комунікацію та потребує уточнення категоріального апарату для цілей нашого дослідження.

%%% PART A: Компетентнісний підхід у педагогіці: витоки та засадничі положення %%%

Для обґрунтування теоретичних засад формування ІКК звернімося до компетентнісного підходу, який утвердився в педагогічній науці як методологічна основа модернізації освітніх систем. Компетентнісний підхід орієнтує освіту не на відтворення знань, а на формування здатності застосовувати знання, уміння, навички й особистісні якості у конкретних професійних та життєвих ситуаціях~\cite{Musiienko2025DesignEducation}.

О.\,І.~Локшина, аналізуючи компетентнісний підхід як концептуальну основу зв'язку між освітою і ринком праці, обґрунтувала розмежування понять <<компетенція>> та <<компетентність>>, де компетенція тлумачиться як соціально задана норма до освітньої підготовки, необхідна для якісної продуктивної діяльності у певній сфері, а компетентність~--- як особистісна якість, що інтегрує знання, уміння, мотивацію, ціннісні орієнтації та досвід діяльності~\cite{Lokshyna2014Competency}. За О.\,І.~Локшиною, компетентність є актуальним, сформованим особистісним утворенням, що проявляється у взаємодії людини зі світом, тобто компетентність завжди має діяльнісний характер і не зводиться до формального знання. Це положення є засадничим для нашого дослідження, оскільки ІКК дизайнера формується не через пасивне засвоєння теоретичних знань про інформаційні технології, а через активну діяльність у цифровому професійному середовищі.

О.\,І.~Пометун, розвиваючи ідеї компетентнісного підходу в українському освітньому контексті, виокремила інформаційну компетентність як одну з ключових освітніх компетентностей, що передбачає уміння самостійно шукати, аналізувати, відбирати потрібну інформацію, організовувати, перетворювати, зберігати та передавати її за допомогою реальних об'єктів та інформаційних технологій~\cite{Pometun2004Competency}. Дослідниця наголошує на тому, що компетенція є соціальною вимогою, наперед заданою нормою підготовки, тоді як компетентність~--- особистісне надбання людини, результат розвитку. Це розмежування дозволяє коректно ставити дидактичні цілі: формування ІКК дизайнера є процесом набуття особистісної компетентності на основі соціально обумовленої компетенції.

У контексті реформування освітніх систем Н.\,М.~Бібік обґрунтувала необхідність компетентнісного підходу для модернізації змісту освіти в Україні, визначивши компетентність як <<результативно-діяльнісну характеристику освіти>>, що відображає спроможність випускника застосовувати набуті знання у нових, нестандартних ситуаціях~\cite{Bibik2004Competency}. Н.\,М.~Бібік підкреслює, що компетентнісний підхід не заперечує значення знань, але переводить акцент зі знань як самоцілі на знання як інструмент розв'язання практичних завдань. Для дизайнерської освіти це положення має особливу вагу, оскільки професійна підготовка дизайнера традиційно будується на інтеграції теоретичних знань і практичної проєктної діяльності.

О.\,В.~Овчарук систематизувала вітчизняний і зарубіжний досвід реалізації компетентнісного підходу, зосередившись на ключових компетентностях в оновленому змісті загальної середньої освіти та обґрунтувавши їх роль у забезпеченні якості освіти~\cite{Ovcharuk2004Competency}. За О.\,В.~Овчарук, ключові компетентності є надпредметними утвореннями, що виходять за межі окремих навчальних дисциплін і забезпечують здатність особистості до соціальної взаємодії, навчання протягом життя та професійної мобільності. Це положення підтверджує правомірність виокремлення ІКК як окремої компетентності: вона не обмежується технічними навичками використання конкретних програмних засобів, а є інтегративним утворенням, що забезпечує ефективну діяльність у цифровому інформаційному середовищі.

М.\,І.~Жалдак, досліджуючи проблеми інформатизації освіти та підготовки вчителів, обґрунтував поняття інформатичної компетентності як здатності орієнтуватися в інформаційному просторі, ефективно використовувати засоби інформаційно-комунікаційних технологій для розв'язання різноманітних завдань професійної та повсякденної діяльності~\cite{Zhaldak2005Competency}. М.\,І.~Жалдак наголошує на нерозривному зв'язку інформаційної та комунікаційної складових: сучасне інформаційне середовище є за своєю природою комунікативним, оскільки інформаційні процеси завжди опосередковуються соціальною взаємодією. Цей висновок є принципово важливим для теоретичного обґрунтування ІКК дизайнера, де комунікативний компонент (візуальна комунікація, презентація проєктних рішень, колаборація у розподілених командах) є невіддільним від інформаційно-технологічного.

Узагальнення позицій зазначених дослідників дозволяє сформулювати засадничі положення для дослідження ІКК:
\begin{enumerate}
    \item компетентність є результатом особистісного розвитку, має діяльнісний характер і не зводиться до знань або вмінь~\cite{Lokshyna2014Competency, Pometun2004Competency};
    \item компетентність формується в діяльності та проявляється у здатності розв'язувати практичні завдання в реальних ситуаціях~\cite{Bibik2004Competency, Ovcharuk2004Competency};
    \item інформаційна та комунікаційна складові утворюють єдність, оскільки інформаційне середовище є за своєю природою комунікативним~\cite{Zhaldak2005Competency};
    \item ключові компетентності є надпредметними утвореннями, що забезпечують інтеграцію знань з різних галузей~\cite{Ovcharuk2004Competency}.
\end{enumerate}

%%% PART B: Термінологічний аналіз суміжних понять %%%

Опираючись на засадничі положення компетентнісного підходу, перейдемо до детального термінологічного аналізу понять, що функціонують у суміжному семантичному полі з ІКК. У сучасному науково-педагогічному дискурсі використовується низка термінів, що частково перетинаються за змістом, проте мають різний обсяг та акценти. Їх чітке розмежування є необхідною передумовою для операціоналізації ІКК у контексті нашого дослідження.

\subsubsection{Інформаційна грамотність (information literacy)}\label{sssec:1.1.1}

Поняття інформаційної грамотності (ІГ) виникло у бібліотечно-інформаційній науці наприкінці ХХ століття і пройшло значну еволюцію від вузького тлумачення як навичок пошуку інформації до широкого розуміння як здатності розпізнавати інформаційну потребу, ефективно шукати, оцінювати та використовувати інформацію. Класичне визначення Американської бібліотечної асоціації (1989) характеризує інформаційно грамотну людину як таку, що здатна <<розпізнати, коли інформація необхідна, і має здатність знайти, оцінити та ефективно використати потрібну інформацію>>~\cite{ALA1989Information}. Надалі це поняття було розширене у межах рамки ACRL (2015), де ІГ визначається як <<набір інтегрованих здатностей, що охоплюють рефлексивне виявлення інформації, розуміння того, як інформація продукується та оцінюється, і використання інформації для створення нових знань та етичної участі у спільнотах навчання>>~\cite{ACRL2015Framework}.

Для дизайнерської професії інформаційна грамотність набуває специфічного змісту: дизайнер має вміти шукати та оцінювати візуальні референси, аналізувати тренди, працювати з базами зображень, патентними базами та портфоліо-платформами. Проте ІГ не охоплює комунікативний вимір (презентація, співпраця, візуальна комунікація) та технологічний вимір (практичне володіння цифровими інструментами проєктування), що обмежує застосовність цього поняття для цілей нашого дослідження.

\subsubsection{ІКТ-компетентність (ICT competence)}\label{sssec:1.1.2}

ІКТ-компетентність (компетентність у сфері інформаційно-комунікаційних технологій) акцентує технологічну складову~--- здатність використовувати апаратне та програмне забезпечення, мережеві технології для розв'язання практичних завдань. Європейська рамка ІКТ-компетентностей (European e-Competence Framework, e-CF) визначає ІКТ-компетентність як <<продемонстровану здатність застосовувати знання, уміння та ставлення для досягнення результатів, що спостерігаються>>, у сфері інформаційно-комунікаційних технологій~\cite{CEN2014eCF}. В українському педагогічному дискурсі М.\,І.~Жалдак визначає ІКТ-компетентність як <<здатність людини орієнтуватися в інформаційному просторі, отримувати інформацію та оперувати нею відповідно до власних потреб і вимог сучасного високотехнологічного інформаційного суспільства>>~\cite{Zhaldak2005Competency}.

Перевагою поняття ІКТ-компетентності є його конкретність: воно безпосередньо пов'язане з інструментальним виміром діяльності. Однак його обмеженням є надмірний акцент на технологічній складовій при недостатній увазі до інформаційно-аналітичного та рефлексивного вимірів. У контексті дизайнерської освіти ІКТ-компетентність описує лише один аспект професійної підготовки~--- володіння цифровими інструментами (Adobe Creative Suite, Figma, 3D-моделювання тощо), не охоплюючи здатність критично оцінювати результати роботи з ШІ, етичні аспекти використання генеративних технологій та рефлексію власного професійного розвитку.

\subsubsection{Цифрова компетентність (digital competence)}\label{sssec:1.1.3}

Цифрова компетентність є найбільш широким поняттям, що набуло інституційного закріплення завдяки Європейській рамці цифрових компетентностей для громадян (DigComp). У рамці DigComp~2.2 цифрова компетентність визначається як <<впевнене, критичне та відповідальне використання і взаємодія з цифровими технологіями для навчання, роботи та участі у суспільстві>>~\cite{Vuorikari2022DigComp}. DigComp виокремлює п'ять вимірів цифрової компетентності: (1)~інформаційна грамотність та грамотність роботи з даними; (2)~комунікація та співпраця; (3)~створення цифрового контенту; (4)~безпека; (5)~розв'язання проблем.

Рамка DigComp~2.2, оновлена у 2022 році, уперше включила навички взаємодії зі штучним інтелектом як складову цифрової компетентності, визнавши необхідність критичного ставлення до ШІ-генерованого контенту та розуміння принципів роботи алгоритмів~\cite{Vuorikari2022DigComp}. Для педагогічної діяльності розроблено окрему рамку~--- DigCompEdu, яка конкретизує цифрову компетентність педагогів~\cite{Redecker2017DigCompEdu}.

Цифрова компетентність є концептуально найближчою до ІКК, проте має універсальний (не професійно-специфічний) характер. Вона описує загальногромадянські навички, тоді як ІКК дизайнера передбачає врахування специфіки дизайнерської діяльності~--- візуальної комунікації, проєктного мислення, роботи зі спеціалізованими інструментами, генеративним ШІ та естетичного оцінювання результатів.

\subsubsection{Медіакомпетентність (media competence)}\label{sssec:1.1.4}

Медіакомпетентність фокусується на здатності особистості критично сприймати, аналізувати, оцінювати та створювати медіаповідомлення у різних формах. У вітчизняній педагогіці медіакомпетентність тісно пов'язана з медіаосвітою та медіаграмотністю. Як зазначає Л.\,А.~Найдьонова, медіакомпетентність виходить за межі суто технологічного виміру і включає критичне мислення щодо медіаконтенту, розуміння механізмів медіавпливу, здатність розрізняти достовірну та маніпулятивну інформацію~\cite{Naidonova2013Media}.

Для дизайнерської професії медіакомпетентність має важливе значення, оскільки дизайнер одночасно є споживачем і виробником медіаконтенту. Проте медіакомпетентність не охоплює повною мірою технологічний вимір дизайнерської діяльності (володіння спеціалізованими інструментами проєктування) та рефлексивний вимір (самооцінка професійного розвитку, етика використання ШІ).

\subsubsection{Інформаційно-комунікативна компетентність (ІКК)}\label{sssec:1.1.5}

Інформаційно-комунікативна компетентність (ІКК) є інтегративним поняттям, що синтезує інформаційний, комунікаційний, технологічний і рефлексивний виміри діяльності. На відміну від ІКТ-компетентності, ІКК не обмежується інструментальним використанням технологій, а охоплює здатність здійснювати цілісну інформаційно-комунікаційну діяльність~--- від пошуку й аналізу інформації через її перетворення та комунікацію до критичної рефлексії результатів. На відміну від цифрової компетентності (DigComp), ІКК може бути операціоналізована для конкретної професійної діяльності, що дозволяє враховувати специфіку відповідної галузі.

У контексті нашого дослідження ІКК визначається як \textit{інтегративна особистісна якість, що характеризує здатність і готовність майбутнього дизайнера до ефективного пошуку, аналізу та використання професійно значущої інформації; цифрової комунікації та візуальної презентації; володіння цифровими інструментами проєктування, включаючи засоби штучного інтелекту; критичної рефлексії власної діяльності в цифровому середовищі та неперервного професійного самовдосконалення}.

Таке визначення враховує засадничі положення компетентнісного підходу (діяльнісний характер, інтегративність, орієнтація на практичний результат) та специфіку дизайнерської професії (візуальна комунікація, проєктна діяльність, робота з ШІ-інструментами).

%%% Порівняльна таблиця %%%

Для систематизації результатів термінологічного аналізу подамо порівняльну характеристику розглянутих понять (табл.~\ref{tab:terminology_comparison}).

\begin{table}[!h]
\centering
\caption{Порівняльна характеристика понять, суміжних з ІКК}
\label{tab:terminology_comparison}
\small
\begin{tabular}{|p{2.6cm}|p{2.4cm}|p{2.5cm}|p{2.4cm}|p{2.5cm}|}
\hline
\textbf{Поняття} & \textbf{Домінантний фокус} & \textbf{Ключова рамка} & \textbf{Обмеження для дизайну} & \textbf{Охоплення ШІ} \\
\hline
Інформаційна грамотність & Пошук, оцінка, використання інформації & ACRL Framework & Не охоплює технологічний та комунікативний виміри & Часткове \\
\hline
ІКТ-компетентність & Технологічне використання ІКТ & e-CF & Надмірний акцент на технологіях & Часткове \\
\hline
Цифрова компетентність & Інтегративне використання цифрових технологій & DigComp 2.2 & Універсальний характер, не враховує специфіку професії & Так (з 2022) \\
\hline
Медіа\-компетентність & Критичне сприйняття та створення медіа & Медіаосвіта & Не охоплює технологічний вимір & Часткове \\
\hline
\textbf{ІКК дизайнера} & \textbf{Інтегративна: інформаційно-аналітичний, комунікативний, технологічний, рефлексивний} & \textbf{Авторська модель} & \textbf{Професійно-специфічна} & \textbf{Системне} \\
\hline
\end{tabular}
\end{table}

Як засвідчує табл.~\ref{tab:terminology_comparison}, кожне з розглянутих понять описує певний аспект взаємодії людини з інформаційним середовищем, проте жодне з них не охоплює всі виміри, необхідні для характеристики інформаційно-комунікаційної діяльності дизайнера. Інформаційна грамотність залишається в межах пошуку й оцінювання інформації; ІКТ-компетентність зосереджена на технологічній складовій; цифрова компетентність є надто загальною; медіакомпетентність не враховує специфіку інструментального виміру дизайнерської діяльності. Саме ІКК як інтегративне поняття, що синтезує всі чотири виміри (інформаційно-аналітичний, комунікативний, технологічний, рефлексивний), є найбільш адекватним для опису цілісної готовності дизайнера до професійної діяльності в цифровому середовищі.

%%% PART C: ІКК у контексті дизайнерської професії %%%

Визначивши місце ІКК серед суміжних понять, розглянемо специфіку інформаційно-комунікативної компетентності у контексті дизайнерської професії. Дизайнерська діяльність має низку характерних особливостей, що зумовлюють специфічний зміст кожного компонента ІКК.

По-перше, дизайн є \textit{візуально-комунікативною діяльністю}. Дизайнер створює візуальні повідомлення, що мають бути зрозумілими цільовій аудиторії, відповідати естетичним та функціональним вимогам, враховувати культурний контекст. Це означає, що комунікативний компонент ІКК дизайнера суттєво відрізняється від комунікативного компонента інших професій: він включає не лише вербальну й текстову, а й візуальну комунікацію, презентацію проєктних рішень, вміння формулювати та захищати дизайн-концепції перед замовником та командою~\cite{Musiienko2025DesignEducation}.

По-друге, дизайнерська діяльність є \textit{інструментально насиченою}. Сучасний дизайнер використовує широкий спектр цифрових інструментів: від графічних редакторів (Adobe Photoshop, Illustrator) та інструментів UX/UI-дизайну (Figma, Sketch) до засобів 3D-моделювання (Blender, 3ds Max), генеративного ШІ (Midjourney, DALL-E, Stable Diffusion) та інструментів автоматизації (скриптування, prompt engineering). Технологічний компонент ІКК дизайнера відповідно охоплює не лише базове використання програмного забезпечення, а й адаптивний вибір інструментів, інтеграцію ШІ-засобів у проєктний робочий процес та prompt engineering як новий вимір технологічної компетентності.

По-третє, дизайнерська діяльність є \textit{проєктною та ітеративною}. Дизайн-процес будується на циклі <<дослідження -- концепція -- прототип -- тестування -- ітерація>>, що потребує постійної рефлексії, самооцінки та коригування діяльності. Рефлексивний компонент ІКК дизайнера включає здатність оцінювати якість власної інформаційно-комунікаційної діяльності, критично аналізувати результати роботи з ШІ-інструментами, усвідомлювати етичні аспекти використання генеративних технологій (авторське право, упередження алгоритмів, автентичність дизайну).

По-четверте, дизайн зазнав радикальної \textit{ШІ-трансформації}. Систематичний огляд 156 досліджень генеративного ШІ в дизайн-освіті засвідчив 6,4-кратне зростання публікацій за останні роки та парадигмальний зсув від <<дизайнера-виробника>> до <<дизайнера-куратора>>, який оцінює, редагує та спрямовує результати, згенеровані ШІ~\cite{Musiienko2026ScopingReview}. Цей зсув змінює саму природу інформаційно-комунікативної діяльності дизайнера: замість безпосереднього створення візуального продукту дизайнер дедалі частіше формулює завдання для ШІ (prompt engineering), оцінює згенеровані варіанти, курує результати та інтегрує їх у цілісний проєкт.

Зазначені особливості дизайнерської професії зумовлюють необхідність операціоналізації ІКК з урахуванням специфіки дизайнерської діяльності. На підставі проведеного аналізу нами обґрунтовано 4-компонентну структуру ІКК дизайнера.

%%% PART D: 4-компонентна структура ІКК дизайнера %%%

Структура ІКК дизайнера включає чотири взаємопов'язані компоненти: інформаційно-аналітичний, комунікативний, технологічний та рефлексивний. Кожен компонент має специфічний зміст, зумовлений особливостями дизайнерської професії в еру штучного інтелекту.

\textbf{1. Інформаційно-аналітичний компонент} охоплює здатність дизайнера до ефективного пошуку, оцінювання, систематизації та використання професійно значущої інформації. Зміст цього компонента включає:

\begin{itemize}
    \item пошук та аналіз візуальних референсів, дизайн-трендів, кейсів та найкращих практик з використанням спеціалізованих платформ (Behance, Dribbble, Pinterest, Dezeen тощо);
    \item критичне оцінювання якості та достовірності інформаційних джерел, включаючи контент, згенерований ШІ~\cite{Musiienko2026ScopingReview};
    \item аналіз цільової аудиторії, користувацьких потреб та контексту використання дизайн-продукту;
    \item систематизація та структурування зібраної інформації для формування дизайн-брифу та концепції;
    \item здатність працювати з великими обсягами візуальної інформації, виокремлювати суттєве, встановлювати зв'язки між різнорідними даними;
    \item розуміння принципів роботи алгоритмів рекомендацій та пошукових систем, які впливають на результати інформаційного пошуку.
\end{itemize}

Інформаційно-аналітичний компонент ІКК дизайнера виходить за межі класичної інформаційної грамотності, оскільки включає специфічні для дизайну навички: аналіз візуальних джерел, робота з портфоліо-платформами, критичне оцінювання ШІ-генерованого контенту. В умовах стрімкого зростання обсягів ШІ-генерованих зображень~\cite{Musiienko2026ScopingReview} здатність відрізняти автентичний дизайн від штучно згенерованого набуває критичного значення.

\textbf{2. Комунікативний компонент} характеризує здатність дизайнера до ефективної професійної комунікації в цифровому середовищі. Зміст цього компонента включає:

\begin{itemize}
    \item візуальну комунікацію~--- здатність створювати зрозумілі, естетично якісні та функціональні візуальні повідомлення для різних цільових аудиторій;
    \item презентацію проєктних рішень~--- уміння переконливо подати дизайн-концепцію замовнику, команді, експертам з використанням цифрових засобів (Figma, Miro, презентаційні платформи);
    \item цифрову колаборацію~--- здатність ефективно працювати у розподілених командах з використанням інструментів спільної роботи (Slack, Notion, Trello, реальночасова колаборація у Figma);
    \item професійне мережування~--- формування та підтримка професійних контактів у цифровому середовищі, участь у професійних спільнотах, формування онлайн-портфоліо;
    \item мультимодальну комунікацію~--- здатність поєднувати текст, зображення, відео, анімацію та інтерактивні елементи для ефективного донесення ідей;
    \item комунікацію з ШІ-системами~--- здатність формулювати чіткі та ефективні запити (prompt engineering) для отримання бажаних результатів від генеративних ШІ-інструментів~\cite{Musiienko2026ScopingReview}.
\end{itemize}

Комунікативний компонент ІКК дизайнера є суттєво ширшим за традиційне розуміння комунікативної компетентності, оскільки включає специфічний для дизайну вимір~--- візуальну комунікацію, а також новий вимір~--- комунікацію з ШІ-системами (prompt engineering), що є формою людино-машинної взаємодії, де якість <<діалогу>> визначає якість результату.

\textbf{3. Технологічний компонент} охоплює здатність дизайнера до практичного використання цифрових інструментів проєктування, включаючи традиційні засоби та засоби штучного інтелекту. Зміст цього компонента включає:

\begin{itemize}
    \item володіння базовими інструментами графічного дизайну (Adobe Creative Suite, Affinity, CorelDRAW);
    \item володіння інструментами UX/UI-дизайну та прототипування (Figma, Sketch, Adobe XD);
    \item базові навички 3D-моделювання та візуалізації (Blender, 3ds Max, SketchUp);
    \item використання генеративних ШІ-інструментів (Midjourney, DALL-E, Stable Diffusion, Adobe Firefly) для створення та модифікації візуального контенту;
    \item prompt engineering~--- уміння формулювати ефективні текстові та візуальні запити для ШІ-систем, включаючи ітеративне вдосконалення промптів~\cite{Musiienko2026ScopingReview};
    \item адаптивний вибір інструментів~--- здатність обирати оптимальний інструмент (традиційний чи ШІ-орієнтований) залежно від завдання, контексту та обмежень;
    \item інтеграція ШІ-засобів у робочий процес~--- здатність поєднувати ШІ-генеровані елементи з авторською роботою у цілісний дизайн-продукт;
    \item базове розуміння принципів роботи ШІ-алгоритмів (генеративні змагальні мережі, дифузійні моделі, трансформери), що уможливлює свідоме та ефективне використання ШІ-інструментів.
\end{itemize}

Технологічний компонент ІКК дизайнера суттєво відрізняється від ІКТ-компетентності загального характеру, оскільки включає специфічні для дизайну інструменти та~--- що є принципово новим~--- навички роботи з генеративним ШІ. Prompt engineering як новий вимір технологічної компетентності стає невід'ємною частиною професійного інструментарію дизайнера~\cite{Musiienko2026ScopingReview}.

\textbf{4. Рефлексивний компонент} характеризує здатність дизайнера до критичного аналізу, самооцінки та саморегуляції інформаційно-комунікативної діяльності. Зміст цього компонента включає:

\begin{itemize}
    \item самооцінку рівня розвитку ІКК~--- здатність адекватно оцінювати власний рівень інформаційно-аналітичних, комунікативних та технологічних навичок, визначати зони розвитку;
    \item критичне оцінювання ШІ-генерованих результатів~--- здатність виявляти артефакти, упередження, стереотипи та етичні проблеми у контенті, згенерованому ШІ~\cite{Musiienko2026ScopingReview};
    \item етику використання ШІ в дизайні~--- усвідомлення питань авторського права, інтелектуальної власності, автентичності дизайну, прозорості використання ШІ-інструментів;
    \item критичне ставлення до алгоритмічних упереджень~--- розуміння того, що ШІ-моделі відтворюють упередження навчальних даних (гендерні, культурні, расові) і здатність компенсувати ці упередження в дизайн-рішеннях;
    \item планування неперервного професійного розвитку~--- здатність визначати потреби у навчанні, обирати форми підвищення кваліфікації, адаптуватися до нових технологій;
    \item рефлексію професійної ідентичності~--- осмислення власної ролі як дизайнера в умовах ШІ-трансформації, визначення балансу між авторським та ШІ-асистованим творенням.
\end{itemize}

Рефлексивний компонент є системоутворювальним у структурі ІКК дизайнера, оскільки забезпечує усвідомленість та саморегуляцію інформаційно-комунікативної діяльності. Без рефлексивного компонента інші компоненти ІКК залишаються набором розрізнених навичок, що не утворюють цілісної компетентності. Саме рефлексивний компонент забезпечує перехід від механічного використання інструментів до осмисленої, етично відповідальної та адаптивної професійної діяльності.

%%% Структурна модель ІКК (TikZ-фігура) %%%

Взаємозв'язок компонентів ІКК дизайнера візуалізовано на рис.~\ref{fig:ikk_model}.

\begin{figure}[!h]
\centering
\resizebox{\textwidth}{!}{%
\begin{tikzpicture}

% Title
\node[dia/header=14cm] (title) {СТРУКТУРА ІКК ДИЗАЙНЕРА};

% Four component boxes — two rows
\node[dia/lrbox=5.8cm, below=1.5cm of title.south, xshift=-3.8cm] (info) {%
  \textbf{1. Інформаційно-аналітичний}\\[3pt]
  \textbullet~пошук та аналіз інформації;\\
  \textbullet~оцінювання джерел (вкл.\ ШІ);\\
  \textbullet~систематизація даних;\\
  \textbullet~формування дизайн-брифу%
};

\node[dia/lrbox=5.8cm, below=1.5cm of title.south, xshift=3.8cm] (comm) {%
  \textbf{2. Комунікативний}\\[3pt]
  \textbullet~візуальна комунікація;\\
  \textbullet~презентація проєктів;\\
  \textbullet~цифрова колаборація;\\
  \textbullet~prompt engineering%
};

\node[dia/lrbox=5.8cm, below=1.2cm of info] (tech) {%
  \textbf{3. Технологічний}\\[3pt]
  \textbullet~цифрові інструменти дизайну;\\
  \textbullet~генеративний ШІ;\\
  \textbullet~адаптивний вибір засобів;\\
  \textbullet~інтеграція ШІ у робочий процес%
};

\node[dia/lrbox=5.8cm, below=1.2cm of comm] (refl) {%
  \textbf{4. Рефлексивний}\\[3pt]
  \textbullet~самооцінка розвитку ІКК;\\
  \textbullet~критичне оцінювання ШІ;\\
  \textbullet~етика використання ШІ;\\
  \textbullet~неперервний розвиток%
};

% Arrows from title to top row
\draw[dia/arrow] (title.south) -- (info.north);
\draw[dia/arrow] (title.south) -- (comm.north);

% Arrows from top row to bottom row
\draw[dia/arrow] (info.south) -- (tech.north);
\draw[dia/arrow] (comm.south) -- (refl.north);

% Bidirectional arrows between components in same row
\draw[dia/darrow] (info.east) -- (comm.west);
\draw[dia/darrow] (tech.east) -- (refl.west);

% Cross arrows
\draw[dia/darrow] (info.south east) -- (refl.north west);
\draw[dia/darrow] (comm.south west) -- (tech.north east);

% Result box
\node[dia/rbox=12cm, below=1.5cm of $(tech.south)!0.5!(refl.south)$] (result) {%
  \textbf{ІКК дизайнера:} інтегративна готовність до професійної діяльності\\
  в цифровому середовищі ери ШІ%
};

% Arrows from bottom row to result
\draw[dia/arrow] (tech.south) -- (result.north);
\draw[dia/arrow] (refl.south) -- (result.north);

\end{tikzpicture}%
}
\caption{Структурна модель інформаційно-комунікативної компетентності дизайнера}
\label{fig:ikk_model}
\end{figure}

Модель (рис.~\ref{fig:ikk_model}) відображає інтегративний характер ІКК: усі чотири компоненти пов'язані двосторонніми зв'язками, що засвідчує їх взаємозумовленість. Інформаційно-аналітичний компонент забезпечує підґрунтя для осмисленого використання технологій (технологічний компонент) та ефективної комунікації (комунікативний компонент). Комунікативний компонент підтримує рефлексивний через механізми зворотного зв'язку від колег, замовників та аудиторії. Технологічний компонент реалізує інформаційно-аналітичні та комунікативні наміри через конкретні інструменти. Рефлексивний компонент є системоутворювальним, оскільки забезпечує усвідомлену саморегуляцію всієї інформаційно-комунікативної діяльності.

%%% PART E: ІКК у контексті ШІ-трансформації дизайну %%%

Запропонована 4-компонентна модель ІКК дизайнера відображає сучасний стан дизайнерської професії, що перебуває у процесі глибокої трансформації під впливом генеративного штучного інтелекту. Аналіз 156 досліджень генеративного ШІ в дизайн-освіті~\cite{Musiienko2026ScopingReview} засвідчив формування нової парадигми~--- перехід від <<дизайнера-виробника>> до <<дизайнера-куратора>>, що безпосередньо впливає на зміст кожного компонента ІКК.

У парадигмі <<дизайнера-виробника>> технологічний компонент ІКК обмежувався володінням графічними редакторами та інструментами проєктування. У парадигмі <<дизайнера-куратора>> технологічний компонент розширюється за рахунок навичок prompt engineering, критичного оцінювання та курування ШІ-генерованого контенту, інтеграції ШІ-засобів у проєктний робочий процес. Аналогічно трансформується інформаційно-аналітичний компонент: дизайнер має не лише шукати інформацію, а й оцінювати достовірність та якість ШІ-генерованих візуальних матеріалів, що потребує нових аналітичних навичок.

Особливого значення набуває рефлексивний компонент, оскільки ШІ-трансформація дизайну породжує низку етичних проблем, виявлених у систематичному огляді~\cite{Musiienko2026ScopingReview}: питання авторського права на ШІ-генерований контент, алгоритмічні упередження (гендерні, культурні, расові стереотипи у навчальних даних), проблема автентичності дизайну та академічної доброчесності. Рефлексивний компонент ІКК забезпечує здатність дизайнера усвідомлювати ці проблеми та приймати етично обґрунтовані рішення.

Водночас ШІ-трансформація не знецінює інші компоненти ІКК, а переозначує їх. Комунікативний компонент збагачується новим виміром~--- комунікацією з ШІ-системами, де якість промпту безпосередньо визначає якість результату. Інформаційно-аналітичний компонент стає критично важливим в умовах зростання обсягів ШІ-генерованого контенту, коли здатність відрізняти достовірну інформацію від згенерованої набуває стратегічного значення.

%%% PART F: Співвіднесення ІКК з DigComp та іншими рамками %%%

Обґрунтовуючи запропоновану модель ІКК, доцільно співвіднести її з існуючими міжнародними рамками цифрових компетентностей. Рамка DigComp~2.2~\cite{Vuorikari2022DigComp} виокремлює п'ять вимірів цифрової компетентності, які частково перетинаються з компонентами ІКК дизайнера:

\begin{enumerate}
    \item \textit{Інформаційна грамотність та грамотність роботи з даними} (DigComp) $\leftrightarrow$ Інформаційно-аналітичний компонент ІКК~--- спільне ядро (пошук, оцінювання інформації), але ІКК включає специфічний для дизайну зміст (візуальні референси, дизайн-тренди, ШІ-генерований контент).
    \item \textit{Комунікація та співпраця} (DigComp) $\leftrightarrow$ Комунікативний компонент ІКК~--- спільне ядро (цифрова комунікація, колаборація), але ІКК розширює за рахунок візуальної комунікації та prompt engineering.
    \item \textit{Створення цифрового контенту} (DigComp) $\leftrightarrow$ Технологічний компонент ІКК~--- найбільший перетин, але ІКК конкретизує для дизайнерських інструментів та генеративного ШІ.
    \item \textit{Безпека} (DigComp) $\leftrightarrow$ частково відображена у рефлексивному компоненті ІКК (етика, авторське право, захист даних).
    \item \textit{Розв'язання проблем} (DigComp) $\leftrightarrow$ є наскрізним виміром ІКК, що реалізується через усі чотири компоненти.
\end{enumerate}

Таким чином, запропонована модель ІКК дизайнера не суперечить рамці DigComp, а конкретизує її для дизайнерської професії з урахуванням ШІ-трансформації, додаючи специфічні для дизайну виміри (візуальна комунікація, prompt engineering, естетичне оцінювання, етика ШІ в дизайні).

%%% PART G: Рівні сформованості ІКК %%%

Для практичного використання запропонованої моделі ІКК у навчальному процесі доцільно визначити рівні сформованості кожного компонента. Опираючись на дескриптори Національної рамки кваліфікацій України та рамки DigComp~2.2~\cite{Vuorikari2022DigComp}, виокремлюємо три рівні сформованості ІКК дизайнера: початковий, достатній та високий (табл.~\ref{tab:ikk_levels}).

\begin{table}[!h]
\centering
\caption{Рівні сформованості компонентів ІКК дизайнера}
\label{tab:ikk_levels}
\small
\begin{tabular}{|p{2.7cm}|p{3.5cm}|p{3.5cm}|p{3.5cm}|}
\hline
\textbf{Компонент} & \textbf{Початковий} & \textbf{Достатній} & \textbf{Високий} \\
\hline
Інформаційно-аналітичний & Виконує пошук за шаблоном; не розрізняє якість джерел & Здійснює цілеспрямований пошук; оцінює якість інформації; працює з кількома типами джерел & Системно аналізує інформаційне поле; критично оцінює ШІ-генерований контент; формує дизайн-бриф на основі аналізу \\
\hline
Комунікативний & Використовує базові засоби комунікації; візуальна презентація нерозвинена & Ефективно комунікує з командою; створює якісні презентації; використовує колаборативні платформи & Майстерно володіє візуальною комунікацією; ефективно використовує prompt engineering; координує розподілені команди \\
\hline
Технологічний & Використовує 1--2 інструменти на базовому рівні; не працює з ШІ & Впевнено володіє основними інструментами; використовує базові ШІ-засоби & Вільно інтегрує традиційні та ШІ-інструменти; адаптивно обирає засоби; ефективний prompt engineering \\
\hline
Рефлексивний & Не здійснює самооцінку; не усвідомлює етичних аспектів ШІ & Оцінює результати власної діяльності; розуміє базові етичні проблеми ШІ & Системно рефлексує; планує саморозвиток; критично аналізує етику ШІ; формує позицію щодо ШІ в дизайні \\
\hline
\end{tabular}
\end{table}

Визначення рівнів сформованості ІКК (табл.~\ref{tab:ikk_levels}) створює підґрунтя для розроблення діагностичного інструментарію та оцінювання ефективності навчальних програм щодо формування ІКК, що буде використано в подальшому дослідженні (розділ~\ref{sec:2}).

%%% PART H: ІКК в контексті стандарту вищої освіти за спеціальністю 022 Дизайн %%%

Для закріплення теоретичних положень у нормативному контексті розглянемо, як компоненти ІКК дизайнера співвідносяться зі стандартом вищої освіти за спеціальністю 022~Дизайн (стандарт В2). Аналіз загальних та фахових компетентностей, визначених стандартом~\cite{Musiienko2025DesignEducation}, засвідчує, що стандарт В2 імпліцитно містить вимоги до формування ІКК, проте не виокремлює ІКК як цілісну компетентність і не враховує ШІ-вимір.

Зокрема, загальна компетентність ЗК-7 (<<Здатність до пошуку, оброблення та аналізу інформації з різних джерел>>) безпосередньо відповідає інформаційно-аналітичному компоненту ІКК, а ЗК-9 (<<Навички використання інформаційних і комунікаційних технологій>>) корелює з технологічним компонентом. Фахова компетентність ФК-14 (<<Здатність застосовувати сучасні інформаційні технології для розв'язання професійних завдань у галузі дизайну>>) також відображає технологічний вимір ІКК. Водночас стандарт не містить компетентностей, пов'язаних з prompt engineering, критичним оцінюванням ШІ-генерованого контенту, етикою використання ШІ та рефлексією професійної діяльності в контексті ШІ-трансформації.

Ця прогалина підтверджується результатами комплексного аналізу 122 українських освітніх програм дизайну, який виявив 0\% покриття компонентів <<Основи ШІ>>, <<Генеративний ШІ>> та <<Комп'ютерний зір>> при загальному рівні інтеграції ШІ-тематики лише 35,7\%~\cite{Musiienko2025DesignEducation}. Це засвідчує розрив між вимогами сучасної дизайнерської практики, де ШІ вже є невід'ємним інструментом, та змістом освітніх програм, що не відображають цю реальність.

Таким чином, запропонована 4-компонентна модель ІКК дизайнера не лише узгоджується з чинним стандартом В2, а й доповнює та розширює його з урахуванням ШІ-трансформації дизайнерської професії, що є необхідною передумовою для модернізації освітніх програм підготовки дизайнерів.

%%% PART I: Огляд вітчизняних досліджень ІКК та суміжних компетентностей %%%

Проблема формування інформаційно-комунікативної компетентності та суміжних компетентностей має значну дослідницьку традицію в українській педагогічній науці, що уможливлює опору на накопичений теоретичний і методичний доробок.

Значний внесок у дослідження інформаційної компетентності педагогів здійснили О.\,М.~Спірін та М.\,І.~Жалдак, які обґрунтували теоретико-методологічні засади формування інформатичної компетентності в системі педагогічної освіти~\cite{Zhaldak2005Competency}. Їхні праці заклали основу для розуміння структури інформаційної компетентності як багатокомпонентного утворення, що інтегрує мотиваційний, когнітивний, операційний та рефлексивний компоненти. Зазначений підхід до структурування компетентності є методологічно близьким до запропонованої нами моделі ІКК дизайнера, де також виокремлено чотири компоненти, проте зміст кожного компонента конкретизовано для дизайнерської професії з урахуванням ШІ-виміру.

К.\,П.~Осадча досліджує проблеми професійної підготовки фахівців у цифровому середовищі, зокрема формування цифрових навичок та використання електронного навчання у підготовці фахівців з дизайну~\cite{Osadcha2024DigitalDesign, Osadcha2024DigitalSkills}. Її дослідження впливу сучасних тенденцій цифрового мистецтва на зміст навчання з комп'ютерної графіки та цифрового дизайну~\cite{Osadcha2021DigitalArt} безпосередньо стосуються технологічного компоненту ІКК дизайнера: зміна інструментарію цифрового дизайну вимагає відповідної трансформації змісту професійної підготовки. Крім того, у нещодавній роботі К.\,П.~Осадча та В.\,В.~Осадчий аналізують практику використання ШІ-інструментів у професійній підготовці педагогів~\cite{Osadcha2025AIEducation}, що засвідчує зростаючу актуальність інтеграції ШІ у вищу освіту.

Проблему формування ІКТ-компетентності у вищій освіті розглядали В.\,Ю.~Биков та М.\,П.~Лещенко, зосереджуючи увагу на цифровій трансформації освіти та необхідності системної модернізації підготовки фахівців у контексті інформаційного суспільства~\cite{Bykov2016DigitalEdu}. Їхні дослідження засвідчують, що формування ІКТ-компетентності не може бути ізольованим від загальної професійної підготовки, а має бути інтегроване у зміст фахових дисциплін. Цей висновок повністю підтверджується нашим аналізом українських освітніх програм дизайну~\cite{Musiienko2025DesignEducation}: успішне формування ІКК дизайнера потребує не окремих курсів з ІКТ, а наскрізної інтеграції цифрових та ШІ-компонентів у проєктні, теоретичні та практичні дисципліни.

Важливим для нашого дослідження є доробок О.\,П.~Пінчук та О.\,М.~Соколюк, які досліджували проблеми використання цифрових засобів в освітньому процесі та формування цифрової компетентності здобувачів освіти~\cite{Pinchuk2022DigitalTools}. Їхні праці актуалізують питання педагогічної доцільності використання цифрових технологій, наголошуючи на тому, що технологія є засобом, а не метою освітнього процесу. Це положення є важливим для формування рефлексивного компонента ІКК дизайнера: студент має не лише опанувати ШІ-інструменти, а й критично оцінювати доцільність їх використання у конкретному проєктному завданні.

У контексті дизайнерської освіти дослідження ІКК є менш систематизованими. Окремі аспекти інформаційної та цифрової компетентності дизайнерів розглядалися у працях, присвячених комп'ютерному дизайну, мультимедійним технологіям в освіті та формуванню проєктної компетентності майбутніх дизайнерів. Однак комплексного дослідження ІКК як інтегративної компетентності, специфічної для дизайнерської професії в контексті ШІ-трансформації, в українській педагогічній науці на момент нашого дослідження не проводилося, що підтверджує його наукову новизну.

%%% PART J: Зарубіжні дослідження %%%

У зарубіжній педагогічній науці проблема цифрових компетентностей дизайнерів набуває дедалі більшої актуальності у зв'язку зі стрімким розвитком генеративного ШІ. Систематичний огляд~\cite{Musiienko2026ScopingReview} виявив 6,4-кратне зростання кількості публікацій з проблематики генеративного ШІ в дизайн-освіті, що засвідчує формування нового дослідницького напряму.

Серед ключових тематичних кластерів, виявлених у систематичному огляді~\cite{Musiienko2026ScopingReview}, найбільшу частку становлять дослідження впливу ШІ на креативність (35,9\%), що безпосередньо пов'язано з рефлексивним компонентом ІКК: дизайнер має критично оцінювати, чи ШІ сприяє креативності, чи обмежує її, замінюючи авторське творення шаблонними рішеннями. Кластер оцінювання (27,6\%) актуалізує питання критеріїв оцінки ШІ-асистованих та ШІ-генерованих дизайн-продуктів. Кластер навчальних програм (22,4\%) безпосередньо стосується проблеми інтеграції ШІ-компонентів у професійну підготовку дизайнерів.

Важливим висновком систематичного огляду є виявлення значної <<атеоретичної прогалини>>: 73,7\% досліджень генеративного ШІ в дизайн-освіті не використовують жодної теоретичної рамки~\cite{Musiienko2026ScopingReview}. Це означає, що інтеграція ШІ в дизайн-освіту здійснюється переважно інтуїтивно, без опори на теоретичні моделі компетентностей, педагогічного дизайну або технологічного прийняття. Запропонована нами 4-компонентна модель ІКК дизайнера покликана частково заповнити цю прогалину, надаючи теоретичну основу для систематичної інтеграції ШІ-компонентів у навчальні програми.

%%% PART K: Роль ІКК у системі професійних компетентностей дизайнера %%%

Визначивши поняття, структуру та специфіку ІКК, розглянемо місце ІКК у загальній системі професійних компетентностей дизайнера. ІКК не є ізольованою компетентністю, а виступає наскрізним утворенням, що забезпечує ефективність реалізації інших професійних компетентностей у цифровому середовищі.

Проєктна компетентність (здатність до повного циклу дизайн-проєктування: від дослідження до реалізації) безпосередньо спирається на ІКК, оскільки кожен етап проєктного процесу потребує інформаційно-комунікаційних дій: пошук та аналіз інформації (інформаційно-аналітичний компонент), комунікація із замовником та командою (комунікативний компонент), використання цифрових інструментів (технологічний компонент), рефлексія результатів (рефлексивний компонент).

Естетична компетентність (здатність до формування та реалізації естетичних рішень) в умовах ШІ-трансформації набуває нового змісту: дизайнер має естетично оцінювати не лише власні рішення, а й ШІ-генерований контент, що потребує розвиненого рефлексивного компонента ІКК.

Дослідницька компетентність (здатність до дослідження користувачів, ринку, технологій) фактично реалізується через інформаційно-аналітичний компонент ІКК з використанням цифрових інструментів дослідження (онлайн-опитування, аналітика, A/B тестування).

Комунікаційна компетентність загального характеру конкретизується у комунікативному компоненті ІКК через специфічні для дизайну форми: візуальну комунікацію, презентацію проєктів, prompt engineering.

Таким чином, ІКК виступає інтегративним, наскрізним компонентом професійної підготовки дизайнера, що забезпечує реалізацію усіх інших професійних компетентностей у цифровому середовищі. Без сформованої ІКК інші компетентності не можуть бути повноцінно реалізовані в сучасних умовах, де цифрові технології та ШІ пронизують усі аспекти дизайнерської діяльності.

%%% PART L: Узагальнення %%%

Проведений аналіз ІКК як педагогічної категорії дозволяє сформулювати такі узагальнення.

Термінологічний аналіз засвідчив, що ІКК є найбільш адекватним поняттям для опису цілісної інформаційно-комунікаційної готовності дизайнера, оскільки інтегрує інформаційно-аналітичний, комунікативний, технологічний та рефлексивний виміри, тоді як суміжні поняття (інформаційна грамотність, ІКТ-компетентність, цифрова компетентність, медіакомпетентність) описують лише окремі аспекти цієї готовності.

Запропоновано визначення ІКК дизайнера як інтегративної особистісної якості, що характеризує здатність і готовність до ефективної інформаційно-комунікаційної діяльності у цифровому професійному середовищі з урахуванням ШІ-трансформації дизайнерської професії.

Обґрунтовано 4-компонентну структуру ІКК дизайнера: (1)~інформаційно-аналітичний компонент (пошук, оцінювання та використання інформації, включаючи ШІ-генерований контент); (2)~комунікативний компонент (візуальна комунікація, презентація, колаборація, prompt engineering); (3)~технологічний компонент (цифрові інструменти дизайну, генеративний ШІ, адаптивний вибір засобів); (4)~рефлексивний компонент (самооцінка, етика ШІ, неперервний розвиток).

Установлено, що ІКК є наскрізною компетентністю, яка забезпечує реалізацію інших професійних компетентностей дизайнера у цифровому середовищі. ШІ-трансформація дизайнерської професії зумовлює необхідність систематичного перегляду змісту освітніх програм з метою інтеграції ШІ-компонентів, що є предметом подальшого аналізу в розділі~\ref{sec:2}.

Для з'ясування масштабів та специфіки ШІ-трансформації дизайн-освіти, що безпосередньо впливає на зміст та актуальність ІКК дизайнера, необхідно звернутися до результатів систематичного огляду сучасних досліджень генеративного ШІ в дизайн-освіті (підрозділ~\ref{subsec:1.2}).
